The method has four stages: (1) construct a semantically embedded paper corpus, (2) build a citation-based DAG, (3) estimate inheritance weights from cited parents to each focal paper, and (4) derive contribution and lineage statistics from the fitted decomposition.

\subsection{Construction of science space}
Let each paper $p_i$ be represented by title and abstract text. We encode this text using an appropriate transformer (e.g.  SciBERT) to obtain an embedding $\mathbf{e}_i \in \mathbb{R}^d$. Embeddings are $\ell_2$-normalized to stabilize cosine geometry and reduce scale effects between documents of different length (large documents having large magnitudes and seen as "more simular").

To improve robustness, we apply light preprocessing (Unicode normalization, citation-marker removal, and deduplication by identifier).

\subsection{Construction of citation DAG}
From citation metadata, we create a directed edge $(j \rightarrow i)$ when paper $p_i$ cites paper $p_j$. For each focal paper, we define its parent set $\mathcal{P}(i)$ as cited predecessors retained after filtering. We then enforce acyclicity by removing chronology-violating edges where publication year$(p_j)$ exceeds publication year$(p_i)$.

\subsection{Inheritance inference}
For each eligible paper $p_i$, we model its embedding as
\begin{equation}
\mathbf{e}_i \approx \sum_{j \in \mathcal{P}(i)} w_{ij}\,\mathbf{e}_j + \mathbf{r}_i,
\end{equation}
where $0 \le w_{ij} \le 1$ and $\mathbf{r}_i$ is a residual contribution vector.

Weights are estimated with non-negative constrained least squares and simplex regularization:
\begin{equation}
\min_{\mathbf{w}_i \ge 0}\; \left\|\mathbf{e}_i-\mathbf{E}_{\mathcal{P}(i)}\mathbf{w}_i\right\|_2^2 + \lambda\|\mathbf{w}_i\|_2^2
\quad\text{s.t.}\quad
\sum_{j \in \mathcal{P}(i)} w_{ij} \le 1,
\end{equation}
with $\mathbf{E}_{\mathcal{P}(i)}$ the parent-embedding matrix.

The non-negativity constraint enforces interpretable additive inheritance, and limiting the total parent weight to at most 1 prevents over-attribution when parent embeddings are highly collinear. The unassigned mass,
\begin{equation}
c_i = 1 - \sum_{j \in \mathcal{P}(i)} w_{ij},
\end{equation}
serves as a scalar \emph{contribution factor} (higher implies less explainability from cited parents). The vector residual
\begin{equation}
\mathbf{r}_i = \mathbf{e}_i-\sum_{j \in \mathcal{P}(i)} w_{ij}\mathbf{e}_j
\end{equation}
captures direction-specific novelty for downstream semantic inspection.

\subsection{Handling non-orthogonal parents}
Cited parents are typically correlated, which can make individual inheritance weights unstable even when overall reconstruction is good. We therefore use ridge regularization ($\lambda>0$) to reduce variance in coefficient allocation and improve numerical stability. We additionally apply bootstrap re-estimation over parent subsets to quantify uncertainty in $w_{ij}$, reporting confidence intervals alongside point estimates to distinguish robust inheritance edges from ill-conditioned assignments.

\subsection{Derived metrics and analyses}
From fitted weights we compute:
\begin{itemize}
    \item \textbf{Edge inheritance strength} $w_{ij}$ for lineage tracing.
    \item \textbf{Paper contribution and contribution  magnitude} $\mathbf{r}_i$, and  $c_i$ or  $\|\mathbf{r}_i\|_2$. % Have some thinking time of what c_i and ||r_i||_2  individually represent. 
    \item \textbf{Lineage concentration} (entropy of $\mathbf{w}_i$), indicating whether inheritance is diffuse or dominated by a few parents.
    \item \textbf{Cross-field transfer} by aggregating inheritance mass across discipline labels.
\end{itemize}

These quantities support both ranking-style comparisons and structural analyses of knowledge evolution in the citation DAG.
